\chapter{Conclusion}
\label{ch:conclusion}

A direct comparison of the co-digestion and glycerol only phases reveals that the co-digestion of food waste and glycerol resulted in higher biogas yields and more stable operation. The pH was decreasing faster with the co-digestion, where with glycerol only phase it wasn't decreasing that fast, and was even growing with the potassium bicarbonate that was added to maintain the pH. These results support the hypothesis that glycerol, when co-digested with food waste, can enhance biogas production by balancing the C/N ratio and providing readily available carbon, whereas glycerol alone may work, but would need a different setup and ajustments.

Using crude glycerol, with some pure glycerol to make the necessary ajustment to balance the C/N ratio, could be a viable solution for the glycerol problem. It's only major problem, is that crude glycerol contains a lot of inpurities that can affect the anaerobic digestion, and for a system that can reliably consume it, it will need a better control of it's substrate.

Future studies could involve on how to use better the crude glycerol as a feed, to try subside it's downside of it's impurities and need for an extra caron rich substrate to balance the feed.
